\section{Related Work}
\label{sec:related}

\textbf{KV recovery after a serving failure.}
D\'ej\`aVu streams KV checkpoints so a failed generative replica can resume
from a recent prefix instead of recomputing it \cite{strati2024dejavu}.
LUMEN coordinates recovery across a distributed serving deployment
\cite{cao2026lumen}. GhostServe keeps a lightweight shadow checkpoint for
fault-tolerant serving \cite{jayakody2026ghostserve}. These systems answer
how to restore a KV cache that has already been selected for protection.
\system asks which Agents and which blocks are worth that copy, using workflow slack and
a Top-$K$ set over the first $L$ layers. Our experiments compare internal
variants of that policy. They are not reproductions of D\'ej\`aVu, LUMEN, or
GhostServe, and they do not claim a lower recovery time than those systems.

\textbf{Live KV reuse.}
PagedAttention and vLLM manage KV as pages so a live worker can share and
schedule it \cite{kwon2023vllm}. SGLang and CachedAttention reuse prefixes
across structured calls and later turns \cite{zheng2024sglang,gao2024cachedattention}.
That reuse disappears when the worker holding the cache fails. A peer copy is
a different object: it must survive the protected failure domain, and it
competes with foreground state for memory. Section~\ref{sec:background}
states the compatibility rules a restored cache still has to satisfy.

\textbf{Agent workflows.}
Parrot treats an LLM application as a program with semantic variables and
schedules dependent calls together \cite{lin2024parrot}. Workflow slack in
\system is the same kind of dependence, used for protection rather than for
batch formation. The block-selection experiment uses single long prefixes, so
it does not exercise the cross-turn history those runtimes are built to
retain.

\textbf{Sparse attention.}
MInference, NOSA, and related native sparse attention select a subset of
historical blocks at inference time \cite{jiang2024minference,huang2026nosav2}.
\system reuses those selection indices as checkpoint counters. It does not
change the kernel's choice. $K$ is the number of blocks kept per layer, and
$L$ is how many early layers receive that copy. Section~\ref{sec:eval-native}
finds that $K=16$ covers 4--8\% of distinct blocks and ties or loses to
recency on that denominator. The online run in Section~\ref{sec:evaluation}
places those blocks and moves workflow completion by at most 0.217\,ms
beyond the DRAM checkpoint.

\textbf{Quota control.}
Additive increase and multiplicative decrease is the classic rule for sharing
a congested resource \cite{chiu1989aimd}. The Adapter uses a guarded form of
that rule to change total peer-HBM capacity only. This evaluation does not
vary that quota; the comparison fixes $K$ and $L$ and asks whether the
selected blocks are the right ones.
